\section{Análise de Algoritmos}%
\begin{frame}[label=main]
	Em computação, o termo \alert{algoritmo} poderia ser descrito como um procedimento 
	bem definido que recebe um valor(es) como \emph{entrada} e produz valor(es)
	como \emph{saída}.
	\vfill\pause
	Exemplos?
	\vfill\pause
	E se houvesse um computador infinitamente rápido com memória infinita?
\end{frame}

\begin{frame}[label=main]
	\framesubtitle{Eficiência}

	Um algoritmo deve:
	\begin{itemize}
		\item funcionar corretamente;
		\item executar o mais rápido possível;
		\item utilizar a memória da melhor forma possível.
	\end{itemize}
	\vfill\pause
	Como comparar algoritmos em função do custo de tempo?
	\begin{itemize}
		\item Computadores diferentes podem funcionar em frequências diferentes.
		\item Compiladores podem otimizar o código antes da execução.
		\item Um algoritmo pode ser escrito de formas diferentes conforme a 
		linguagem de programação utilizada.
	\end{itemize}
\end{frame}

\begin{frame}[fragile,label=graduacao]
	\framesubtitle{Exemplo}

	Diferentes algoritmos podem realizar a mesma tarefa com instruções diferentes.
	\vfill
	Problema: ordenar uma lista de números
	\vfill\pause
	\begin{columns}[T]
		\column{.42\textwidth}
			\begin{lstlisting}[morekeywords={menor,insereFinal,remove}]
				listaOrdenada := emptyset
				enquanto numeros != emptyset
				    n := menor(numeros)
				    insereFinal(listaOrdenada,n)
				    remove(numeros,n)
			\end{lstlisting}
		\column{.45\textwidth}
			\begin{lstlisting}[morekeywords={maior,insereInicio,remove}]
				listaOrdenada := emptyset
				enquanto numeros != emptyset
				    n := maior(numeros)
				    insereInicio(listaOrdenada,n)
				    remove(numeros,n)
			\end{lstlisting}
	\end{columns}
\end{frame}


\begin{frame}[label=main]
	\framesubtitle{Contando Operações}

	Pode-se medir ao custo de tempo contando quantas operações são realizadas pelo algoritmo
	\\\vfill
	Mas o que é uma operação?
	\begin{itemize}
		\item Atribuições
		\item Comparações
		\item Instruções de retorno
		\item etc.
	\end{itemize}
	\vfill\pause
	Cada operação demora o mesmo tanto?
\end{frame}

\newcommand{\comentario}[2][1-]{\uncover<#1>{\textcolor{gray}{// #2}}}%
\begin{frame}[fragile,label=main]
	\framesubtitle{Busca por um elemento}

	\uncover<2->{Pior Caso: não existe no vetor}
	\begin{center}
	\begin{minipage}{.58\textwidth}
	\begin{lstlisting}[morekeywords={busca_linear,tamanho}] 
		busca_linear(vetor, chave)
		    i := 0                         @\comentario[2]{1}@
		    enquanto i < tamanho(vetor)    @\comentario[2]{n+1}@
		        se vetor[i] = chave        @\comentario[2]{n}@
		            retorna i
		        i := i + 1                 @\comentario[2]{n}@
		    retorna -1                     @\comentario[2]{1}@
	\end{lstlisting}
	\hfill\pause Total: 3n+3
	\end{minipage}
	\end{center}
	\vfill\pause
	Melhor caso?
\end{frame}

\begin{frame}[fragile,label=graduacao]
	Diferentes algoritmos podem realizar a mesma tarefa com instruções diferentes 
	a um custo maior/menor
	\\ \vfill
	Problema: ordenar os números da mega-sena
	\vfill\pause
	\begin{columns}[T]
		\column{.01\textwidth}
		\column{.4\textwidth}
			\begin{lstlisting}[morekeywords={menor,adicionar,remover}]
				for(int i = 1; i < 60; ++i) {
				  i = menor(numeros);
				  adicionar(lista_ord, i);
				  remover(numeros, i);
				}
			\end{lstlisting}
		\column{.01\textwidth}
			\pause
		\column{.5\textwidth}
			\begin{lstlisting}[morekeywords={em_ordem,embaralhar}]
				lista_ord = numeros;
				while(!em_ordem(lista_ord))
				  lista_ord = embaralhar(numeros);
			\end{lstlisting}
	\end{columns}
\end{frame}

\begin{frame}[label=graduacao]
	As vezes ser ``melhor'' não quer dizer que seja mais eficiente...
	\pause\\\vfill
	John Richard Boyd
	\begin{itemize}
		\item EUA, $^{\star}$1927 $^{\dagger}$1997
		\item piloto, consultor do Pentágono
		\item OODA - tomada de decisões \pause
			\begin{itemize}
				\item Observe
				\item Oriente
				\item Decida
				\item Aja
			\end{itemize}
	\end{itemize}
	
	\pause%\\\vfill
	\begin{beamerboxesrounded}{}
		A \emph{velocidade} de iteração \uncover<3->{pode} supera\uncover<3->{r} 
		a \emph{qualidade} da iteração
	\end{beamerboxesrounded}

	Exemplos:\pause
	\begin{columns}[T]
		\column{.3\textwidth}
			\subitem{Android}
		\column{.3\textwidth}
			\subitem{Google Chrome}
		\column{.3\textwidth}
			\subitem{Firefox}
	\end{columns}
\end{frame}